\section{Antecedentes del problema}

El crecimiento de una empresa viene acompañado naturalmente de un cierto grado de desorden operativo: a más proyectos, más empleados y más desafíos logísticos. A medida que una organización crece, enfrenta distintas dificultades para integrar al personal nuevo de manera eficiente. Los procesos de inducción (esenciales para asegurar que los nuevos empleados se adapten rápidamente y empiecen a contribuir de manera efectiva), suelen volverse cada vez más complejos y exigentes con dicho crecimiento (\cite{Achury2016induccion}). La automatización de procesos de capacitación y entrega de conocimiento mediante Inteligencia Artificial Generativa cuenta con antecedentes prometedores en el ámbito educativo. \textcite{Zamfirescu2023_GAIED} concluyen que la implementación de sistemas de automatización con base en Modelos Preentrenados (GPT, \textit{Generative Pre-trained Models}) ajustados mediante \textit{fine-tuning} ofrece una solución prometedora para mejorar la eficiencia de procesos de capacitación. \textcite{Schmucker2023_GAIED}, mediante un Sistema Tutor Conversacional que emplea dos LLM, evidencian mejoras en la entrega consistente e instantánea de materiales de capacitación, eliminando la necesidad de una supervisión intensiva por parte de un instructor humano. 

El paradigma de Generación Aumentada por Recuperación (RAG) fue formalmente introducido por \textcite{Lewis2020rag}, quienes propusieron acoplar un modelo generador paramétrico con un mecanismo de recuperación no paramétrico sobre un índice denso de documentos, permitiendo anclar las respuestas del modelo en evidencia externa recuperable en tiempo de inferencia. En paralelo, \textcite{Karpukhin2020_DPR} formalizaron la Recuperación Densa de Pasajes (DPR, \textit{Dense Passage Retrieval}) como alternativa a la recuperación léxica clásica, y \textcite{Guu2020_REALM} propusieron REALM, que integra la recuperación directamente en el preentrenamiento del modelo de lenguaje. Estos tres trabajos, publicados entre 2020 y 2021, constituyen lo que en la literatura reciente se reconoce como la fase de \textquotedbl{}RAG fundacional\textquotedbl{}.

Dualbiz SRL, una empresa boliviana de desarrollo de \textit{software} con 18 años de presencia en el mercado local. En constante crecimiento, enfrenta dificultades, tales como: sus procesos de capacitación y entrenamiento del personal nuevo en las áreas de soporte y desarrollo suelen extenderse más allá de lo previsto, generando retrasos en la ejecución de proyectos, un costo mayor al previsto y un alto costo de oportunidad al posponer la incorporación de nuevos proyectos mientras se resuelven los ya comprometidos. El análisis apriori permite identificar que la carencia real de Dualbiz SRL no es la ausencia de contenido de capacitación personalizado, sino la falta de un mecanismo confiable mediante el cual el personal, tanto nuevo como existente pueda consultar y obtener respuestas fundamentadas sobre el conocimiento operativo ya existente en la organización (procesos internos, decisiones técnicas documentadas, historial de incidencias de soporte), sin exponer información confidencial a servicios de terceros ni requerir una inversión en infraestructura de cómputo que resulta inaccesible para una PyME tecnológica. Esta limitación práctica es consistente con la reportada, de forma general, por organizaciones de otros sectores (manufactura, finanzas, salud, servicios profesionales) que gestionan grandes volúmenes de conocimiento interno confidencial sin soluciones de IA accesibles que operen dentro de sus restricciones de privacidad y \textit{hardware}.

De la confluencia de ambos antecedentes emerge la esencia del vacío que motiva esta investigación: existe evidencia teórica consolidada de que RAG y la autorreflexión, a partir de la implementación de modelos críticos, en combinación, mejoran la confiabilidad de los modelos de lenguaje de consulta. Sin embargo, no existe una arquitectura validada que integre ambos mecanismos, ni evidencia empírica de su aplicabilidad contextualizada para el soporte de conocimiento operativo en la industria del \textit{software} como lo es Dualbiz SRL con restricciones de acceso a la informacion y de hardware. \textcite{jeong_improving_2024} reportan que la aplicación de estos métodos a problemas de dominio específico revela una generalización deficiente que conduce a recuperar documentos incorrectos y a emitir juicios inexactos. Ademas, demuestran que los componentes específicos de dominio (recuperador, corpus documental e instrucciones) resultan necesarios para responder adecuadamente para cada dominio. 
